\documentclass[12pt,a4paper]{article}

\usepackage[utf8]{inputenc}
\usepackage[T1]{fontenc}
\usepackage[provide=*,spanish]{babel}
\usepackage{amsmath,amsfonts,amssymb}
\usepackage{geometry}
\usepackage{hyperref}
\usepackage{enumitem}
\usepackage{microtype}
\usepackage{mathpazo}
\usepackage{graphicx}
\usepackage{float}
\usepackage{booktabs}
\usepackage{array}

\geometry{margin=2.5cm}

\hypersetup{
    colorlinks=true,
    linkcolor=blue,
    urlcolor=blue,
    citecolor=blue
}

\setlist[itemize]{itemsep=2pt,topsep=4pt}
\setlist[enumerate]{itemsep=4pt,topsep=4pt}

\usepackage{xcolor}
\usepackage{listings}
\lstset{basicstyle=\ttfamily\footnotesize,breaklines=true,columns=fullflexible,frame=single,rulecolor=\color{gray},showstringspaces=false}
\begin{document}

%==================================================
% CARÁTULA
%==================================================

\hypersetup{pageanchor=false}

\begin{titlepage}
    \centering
    \vspace*{2cm}

    {\LARGE\textbf{Universidad del Valle de Guatemala}}\\[0.5cm]
    {\Large Facultad de Ingeniería}\\[1.5cm]

    {\Huge\textbf{Reporte Escrito}}\\[0.5cm]
    {\Large\textbf{MLflow en la plataforma Databricks}}\\[0.4cm]
    {\Large\textbf{Machine Learning Engineering}}\\[2cm]

    \textbf{Integrantes:}\\[0.3cm]
    Andy Fuentes --- 22944\\
    Diederich Solís --- 22952\\
    Diego Linares --- 221256\\
    Christian Echeverria\\[2cm]

    \textbf{Fecha: octubre de 2026}

    \vfill
\end{titlepage}

\hypersetup{pageanchor=true}
\pagenumbering{arabic}

%==================================================

\section{Introducción y alcance}
Este informe investiga el uso de MLflow en Databricks a partir del pipeline de Champions League del Taller 3 y la Actividad 9, publicado en el repositorio original de Diego Linares [6]. Se analiza su relación con CRISP-DM y la integración del pipeline mediante un notebook y un archivo Python.

MLflow permite organizar experimentos y empaquetar modelos; Databricks integra estas capacidades con su plataforma de datos. El registro gobernado en Unity Catalog y la publicación mediante Model Serving permiten continuar el ciclo de vida después del entrenamiento [1, 3].

\subsection{¿Qué usos puede tener para proyectos de Machine Learning?}
\begin{itemize}
\item \textbf{Comparación de experimentos:} conservar hiperparámetros y resultados para elegir una configuración con evidencia [2].
\item \textbf{Trazabilidad:} relacionar cada ejecución con métricas, archivos y contexto del entrenamiento [2].
\item \textbf{Reutilización:} guardar y cargar un pipeline completo, incluyendo transformaciones y clasificador [5].
\item \textbf{Colaboración:} centralizar versiones y controlar su acceso mediante Unity Catalog [3].
\item \textbf{Operación:} integrar modelos con servicios de inferencia; publicar un modelo requiere configurar el servicio adicional [1].
\end{itemize}

\subsection{Implementación de referencia: Taller 3}
La base es el trabajo \textit{MLflow con scikit-learn} proporcionado por el equipo [6]. Utiliza el dataset de partidos de Champions League y compara SVC, regresión logística, HistGradientBoosting y Random Forest. El repositorio ya contiene un notebook de integración, módulos Python, seguimiento de datos y registro de modelos.

Por ello, el objetivo de esta entrega es trasladar esa implementación a Databricks: conservar el pipeline y cambiar la configuración de seguimiento, las rutas del experimento y, si se utiliza, el registro gobernado en Unity Catalog. La integración reutiliza el preprocesamiento, los clasificadores y las funciones de seguimiento desarrolladas en el taller.

\subsection{Usos en este proyecto}
Se pueden comparar familias por F1 macro, identificar los datos usados por cada run y recuperar el modelo seleccionado por versión o alias. La misma base permite documentar cambios del preprocesamiento y repetir búsquedas de hiperparámetros. Las variables incluyen estadísticas observadas durante el partido. Por ello, esta evaluación clasifica resultados con datos del encuentro; no demuestra capacidad de pronóstico antes del partido.

\newpage
\section{¿Qué etapas de CRISP-DM puede cubrir?}
La siguiente relación es un análisis metodológico del equipo. MLflow apoya especialmente modelado y evaluación; junto con Databricks también participa en preparación y despliegue. El criterio del negocio y la interpretación de datos siguen siendo responsabilidad del equipo.

\renewcommand{\arraystretch}{1.38}
\noindent\begin{tabular}{@{}>{\raggedright\arraybackslash}p{0.22\textwidth} >{\raggedright\arraybackslash}p{0.55\textwidth} >{\raggedright\arraybackslash}p{0.16\textwidth}@{}}
\toprule
\textbf{Etapa} & \textbf{Aplicación concreta} & \textbf{Cobertura}\\
\midrule
Comprensión del negocio & Documentar objetivo, supuestos y criterio de éxito en notas o etiquetas. La herramienta no define la necesidad del negocio. & Indirecta\\
Comprensión de los datos & Conservar referencias y archivos de diagnóstico. El análisis exploratorio se realiza con notebooks y herramientas de datos. & Parcial\\
Preparación de los datos & Guardar el preprocesamiento dentro del pipeline. Feature Store añade reutilización de variables en la plataforma. & Parcial\\
Modelado & Registrar cada configuración, entrenar modelos y conservar el pipeline y sus dependencias. & Directa\\
Evaluación & Comparar métricas de validación y conservar resultados del conjunto de prueba; revisar también criterios de negocio. & Directa\\
Despliegue & Versionar el modelo y cargarlo para inferencia. Model Serving y Jobs requieren configuración adicional en Databricks. & Con integración\\
\bottomrule
\end{tabular}

\subsection{Correspondencia con Champions League}
\begin{enumerate}
\item \textbf{Preparación:} eliminar filas vacías y columnas que revelan el resultado; convertir porcentajes y razones; imputar, escalar y codificar categorías dentro del pipeline.
\item \textbf{Modelado:} ajustar las familias mediante RandomizedSearchCV, con F1 macro y validación estratificada.
\item \textbf{Evaluación:} comparar validación y prueba con la referencia de clase mayoritaria; conservar la matriz de confusión.
\item \textbf{Despliegue:} el taller registra la SVC con el alias \texttt{campeon}; en Databricks debe configurarse el registro correspondiente.
\end{enumerate}

\newpage
\section{Adaptación del pipeline de Champions League}
\subsection{Estructura de la implementación}
El notebook \texttt{03\_pipeline\_sklearn\_mlflow.ipynb} integra el pipeline del taller. El archivo \texttt{seguimiento.py} contiene \texttt{entrenar\_familia}, que registra parámetros, métricas, datasets, matriz de confusión, resultados de búsqueda y el modelo completo [6].

El pipeline usa \texttt{ColumnTransformer} para variables numéricas, porcentajes, razones y categorías; aplica imputación, escalado y One-Hot Encoding, y después \texttt{SelectKBest} y el clasificador. Conserva los transformadores propios del paquete.

\subsection{Resultados reportados por el equipo}
\renewcommand{\arraystretch}{1.2}
\begin{tabular}{lrrr}
\toprule
Familia & F1 macro CV & Exactitud prueba & F1 macro prueba\\
\midrule
SVC & 0.638 & 0.655 & 0.601\\
Regresión logística & 0.611 & 0.621 & 0.542\\
HistGradientBoosting & 0.600 & 0.621 & 0.579\\
Random Forest & 0.583 & 0.621 & 0.576\\
\bottomrule
\end{tabular}

\noindent\textit{Fuente: Taller 3 original, registro manual, página 10 [6]. Estas cifras pertenecen al trabajo original; no son resultados de esta adaptación ni de una ejecución en Databricks.}

\subsection{Cambio para ejecutar en Databricks}
La integración utiliza el servidor de seguimiento administrado de Databricks y una ruta de experimento dentro del workspace:
\begin{lstlisting}[language=Python]
from taller3_pipeline.seguimiento import (
    configurar_mlflow, entrenar_familia)

# La ruta del experimento se configura por usuario.
configurar_mlflow(
    experimento="/Users/usuario/Champions_Resultado",
    uri="databricks")
resultado = entrenar_familia(
    "svc", n_iter=40, folds=5, n_jobs=1)
\end{lstlisting}
El notebook de entrega importa la función del archivo Python de la raíz; ambos reutilizan el mismo paquete instalado y su CSV incluido. La lectura del dataset no depende del directorio de trabajo. La URI de seguimiento es configurable para validar localmente con SQLite y ejecutar después con el servicio administrado de Databricks. El registro opcional en Unity Catalog requiere un nombre \texttt{catalogo.esquema.modelo} y permisos [3]. El alcance de la implementación comprende Experiments, Tracking, Dataset Tracking y MLflow Models. Se guarda el pipeline con su firma, ejemplo de entrada y código de los transformadores, y se comprueba su recarga calculando nuevamente el F1 macro de prueba. La implementación se entrega preparada para Databricks; no se ha ejecutado ni validado en ese servicio. El registro de una versión en Unity Catalog es opcional y no equivale a desplegar un endpoint.

\newpage
\section{Conclusiones}
MLflow permite conservar la relación entre datos, hiperparámetros, métricas y modelos. En el proyecto de Champions League facilita comparar familias con F1 macro y recuperar el pipeline completo, incluyendo sus transformaciones.

La integración con Databricks centraliza los experimentos en un workspace y utiliza el mismo código de entrenamiento desde un notebook o un archivo Python. El registro manual permite conservar la métrica de validación, los resultados de prueba y los artefactos de evaluación de forma explícita.

MLflow ofrece apoyo directo en modelado y evaluación dentro de CRISP-DM. Las herramientas de datos y gobierno de Databricks amplían su participación en preparación y despliegue. El valor de la solución consiste en mantener un proceso trazable y reproducible, con una selección de modelos basada en evidencia.

\section{Repositorio y código del proyecto}
\noindent\textbf{Repositorio de esta entrega:}\\
\url{https://github.com/Andyfer004/Ejercicio5-MLOPS}

\noindent\textbf{Notebook de entrega:}\\
\href{https://github.com/Andyfer004/Ejercicio5-MLOPS/blob/main/notebooks/Champions_Databricks.ipynb}{notebooks/Champions\_Databricks.ipynb}

\noindent\textbf{Archivo Python:}\\
\href{https://github.com/Andyfer004/Ejercicio5-MLOPS/blob/main/pipeline_champions_databricks.py}{pipeline\_champions\_databricks.py}

\noindent\textbf{Paquete y notebook de referencia:}\\
\texttt{src/taller3\_pipeline/} y
\texttt{notebooks/03\_pipeline\_sklearn\_mlflow.ipynb}.

\noindent\textbf{Notebook en Databricks:} el enlace se incorporará después de ejecutar ambos formatos y compartir el notebook con el docente. No hay un enlace ni resultados de Databricks disponibles en esta preparación.

\noindent\textbf{Procedencia y créditos:} pipeline y dataset reutilizados del Taller 3, publicado por Diego Linares, con autoría de Diego Linares, Andy Fuentes, Diederich Solís y Christian Echeverria. Versión de referencia: \texttt{bd619c3cea3e}. El README detalla la instalación y ejecución de ambos formatos.

\section{Referencias oficiales}
\footnotesize
\begin{enumerate}[label={[\arabic*]},itemsep=0pt,topsep=2pt]
\item Databricks. \textit{MLflow on Databricks}.\\\url{https://docs.databricks.com/aws/en/mlflow/}
\item Databricks. \textit{Track model development using MLflow}.\\\url{https://docs.databricks.com/aws/en/mlflow/tracking}
\item Databricks. \textit{Manage model lifecycle in Unity Catalog}.\\\url{https://docs.databricks.com/aws/en/machine-learning/manage-model-lifecycle/}
\item Databricks. \textit{Databricks Feature Store}.\\\url{https://docs.databricks.com/aws/en/machine-learning/feature-store/}
\item MLflow. \textit{mlflow.sklearn: Python API}.\\\url{https://mlflow.org/docs/latest/api_reference/python_api/mlflow.sklearn.html}
\item Equipo del curso. \textit{Taller 3: MLflow con scikit-learn}. Octubre de 2026. PDF original conservado en referencias/Taller3\_MLOPS.pdf y repositorio original: \url{https://github.com/DiegoLinares11/MLFlow}
\end{enumerate}
\noindent Consulta: 16 de octubre de 2026.
\end{document}
